\documentclass[11pt,letterpaper]{article}
\usepackage[margin=1in]{geometry}
\usepackage{amsmath,amssymb,graphicx,booktabs,float}
\usepackage[hidelinks]{hyperref}
\graphicspath{{results/report_assets/}}
\hypersetup{pdftitle={Music Graph Analysis},pdfauthor={}}
\setlength{\emergencystretch}{1em}
\title{Music Graph Analysis}
\author{}
\date{}

\begin{document}
\maketitle
\vspace{-3em}

\section*{E1. Dataset and feature extraction}

The dataset consists of 2,000 tracks by 673 artists from the MTG-Jamendo mood/theme
subset \cite{jamendo,repository}. The aim is to study acoustic similarity,
connectivity and spectral groups, and relate those groups to music tags.
Using seed 521, 12 of 100 shuffled feature archives were selected, and 2,000
eligible tracks were sampled from those archives. This is a two-stage pilot
sample, not a uniform sample of the entire collection.

Each track is represented by 100 features from the supplied
Essentia/AcousticBrainz statistics:
\begin{center}
\begin{tabular}{llr}
\toprule
Group & Features & Dimension\\
\midrule
Timbre & 13 MFCC means and 13 standard deviations & 26\\
Pitch class & 36 HPCP means and 36 standard deviations & 72\\
Rhythm & BPM and onset rate & 2\\
\bottomrule
\end{tabular}
\end{center}
MFCC denotes mel-frequency cepstral coefficients; HPCP denotes harmonic
pitch-class profile. Standard deviations are the square roots of the MFCC
covariance diagonal and HPCP variances.

All features are standardized using population standard deviations. A group of
size $m_g$ is scaled by $1/\sqrt{m_g}$, giving each group equal total variance
before PCA. With $\mu_j$ and $s_j$ denoting feature means and standard deviations,
\[
 \widetilde{x}_{ij}=\frac{x_{ij}-\mu_j}{s_j\sqrt{m_{g(j)}}},
 \qquad f_i=(\widetilde{x}_i-\overline{\widetilde{x}})V_{32}.
\]
Full-SVD PCA retains 32 components without whitening, preserving 94.42\% of the
group-weighted variance. The retained dimension is a hyperparameter.
Genre and other tags are excluded from the feature vectors.

\section*{E2. Pairwise interaction matrix}

For the PCA features $f_i\in\mathbb{R}^{32}$, define Euclidean distances and Gaussian
similarities by
\[
 R_{ij}=\lVert f_i-f_j\rVert_2,\qquad
 W_{ij}=\exp\!\left(-\frac{R_{ij}^2}{2\sigma^2}\right)\ (i\ne j),
 \qquad W_{ii}=0.
\]
The bandwidth $\sigma=1.24392285$ is the median tenth-nearest-neighbor distance,
held fixed for all graphs. The matrices are symmetric and $2000\times2000$,
covering 1,999,000 unordered pairs. The median distance is 2.253910 and the median
off-diagonal similarity is 0.193678. Larger weights indicate greater acoustic
similarity. Three distinct track pairs have identical feature vectors and
similarity one; their track identifiers remain separate.

\section*{E3. kNN graphs and connected components}

For every integer $k=2,\ldots,64$, let $N_k(i)$ contain the $k$ nearest non-self
neighbors of vertex $i$. Distance ties are resolved by node index. The weighted,
undirected union-kNN adjacency is
\[
 A^{(k)}_{ij}=W_{ij}\,
 \mathbf{1}\{j\in N_k(i)\ \text{or}\ i\in N_k(j)\},
 \qquad A^{(k)}_{ii}=0.
\]
Thus an edge is kept if either endpoint selects the other; a vertex's degree may
exceed $k$. Because the Gaussian kernel preserves distance ordering, this is
equivalent to retaining the largest non-self similarities before symmetrization.

\begin{figure}[H]
\centering
\includegraphics[width=.92\linewidth]{latex_connectivity.pdf}
\caption{Connected-component count for all $k=2,\ldots,64$.}
\label{fig:connectivity}
\end{figure}

At $k=2$, the graph has 3,311 edges and components of sizes 1,997 and 3.
All graphs with $k\ge3$ are connected, so the minimum connectivity threshold is
$k_c=3$. Fixed neighbor ordering gives nested edge sets as $k$ increases.
The 1,997-vertex component, containing 99.85\% of the sample, is the large
component analyzed in E4 and E5.

\section*{E4. Degree and local clustering distributions}

Combinatorial statistics use the binary topology $B_{ij}=\mathbf{1}\{A_{ij}>0\}$.
For degree $d_i$ and the number $T_i$ of triangles incident on vertex $i$,
\[
 d_i=\sum_jB_{ij},\qquad C_i=\frac{2T_i}{d_i(d_i-1)}.
\]
Every vertex has degree at least two, so all denominators are nonzero.
Figure~\ref{fig:distributions} shows the empirical probability masses.

\begin{figure}[H]
\centering
\includegraphics[width=\linewidth]{latex_distributions.pdf}
\caption{Degree and local clustering distributions of the large $k=2$ component.
Clustering values are grouped by their exact rational values.}
\label{fig:distributions}
\end{figure}

\begin{center}
\begin{tabular}{lr@{\qquad}lr}
\toprule
Statistic & Value & Statistic & Value\\
\midrule
Vertices / edges & 1,997 / 3,308 & Mean $C_i$ & 0.140104\\
Degree range & 2--19 & Median $C_i$ & 0\\
Mean / median degree & 3.312969 / 3 & Fraction with $C_i=0$ & 65.20\%\\
Degree mode & 2 & Fraction with $C_i=1$ & 8.26\%\\
\bottomrule
\end{tabular}
\end{center}

The degree distribution is discrete and right-skewed, with a dominant mode at
2: 44.32\% of vertices have degree 2 and 81.97\% have degree at most 4.
The observed tail alone does not establish a power law or other parametric family.
The clustering distribution has a large point mass at zero and several
rational-valued spikes. Such spikes are natural at low degrees: when $d_i=2$,
only $C_i=0$ and $C_i=1$ are possible.
The three-vertex component is a triangle, whose degree and clustering
distributions are point masses at 2 and 1, respectively.

\section*{E5. Laplacian spectral analysis}

\subsection*{Normalized Laplacian and the Fiedler value}

Spectral analysis uses the original Gaussian weights. Let
$D_{ii}=\sum_j A_{ij}$ be weighted degree (strength). The symmetric normalized
Laplacian is
\[
 \widehat L=I-D^{-1/2}AD^{-1/2}.
\]
Within a connected component,
$0=\lambda_1<\lambda_2\le\cdots\le\lambda_n\le2$, with zero eigenvector
proportional to $(\sqrt{D_{11}},\ldots,\sqrt{D_{nn}})$.

\begin{figure}[H]
\centering
\includegraphics[width=.90\linewidth]{latex_fiedler.pdf}
\caption{Full-graph Fiedler value. The graph becomes connected at $k_c=3$.}
\label{fig:fiedler}
\end{figure}

For the full 2,000-vertex graph, $\lambda_2(k)$ increases at every step from
0.03215373 at $k=3$ to 0.13746997 at $k=64$, indicating stronger normalized
spectral connectivity. This is an empirical observation, not a general
monotonicity guarantee: adding edges changes both $A$ and $D$.
At $k=2$, the full graph has $\lambda_2=0$, whereas its large connected
component has $\lambda_2=0.01816991$.

\subsection*{Low spectrum and a separated simple eigenvalue}

Figure~\ref{fig:spectrum} shows the 32 smallest eigenvalues of the large
component. To select a separated nonzero eigenvalue, maximize
\[
 g_j=\min\{\lambda_j-\lambda_{j-1},\lambda_{j+1}-\lambda_j\},
 \qquad j=2,\ldots,32.
\]
The 33rd eigenvalue supplies the last candidate's right-hand gap.

\begin{figure}[H]
\centering
\includegraphics[width=\linewidth]{latex_spectrum.pdf}
\caption{The 32 lowest eigenvalues and two-sided isolation gaps.}
\label{fig:spectrum}
\end{figure}

The selected eigenvalue is $\lambda_2=0.0181699110$. Its left and right gaps are
0.0181699110 and 0.0064237727, respectively; the smaller gap is 35.35\% of
$\lambda_2$. The largest eigenpair residual is $1.12\times10^{-15}$, far below
these gaps, supporting numerical simplicity and separation.
The small weighted triangle has eigenvalues $0$, $1.47716388$, and $1.52283612$.

\subsection*{Three-dimensional embedding and nodal domains}

Use the coordinates $(q_2(i),q_3(i),q_4(i))$ from orthonormal eigenvectors.
Their eigenvalues are $0.01816991$, $0.02459368$, and $0.02853374$, respectively.
Writing $Q=[q_2\ q_3\ q_4]$ gives
\[
 \widehat LQ=Q\,\operatorname{diag}(\lambda_2,\lambda_3,\lambda_4),
\]
so these vectors span a three-dimensional invariant subspace.
A strong nodal domain is a connected component of the subgraph induced by
$q_2>0$ or $q_2<0$. Both sign-induced subgraphs are connected, giving domains
of 963 and 1,034 vertices. No eigenvector entries are numerically zero.

\begin{figure}[H]
\centering
\includegraphics[width=\linewidth]{latex_embedding.pdf}
\caption{Two views of the embedding in $(q_2,q_3,q_4)$, colored by the nodal
domains of $q_2$: domain 1 (blue, 963 vertices) and domain 2 (orange, 1,034 vertices).}
\label{fig:embedding}
\end{figure}

The domains are joined by 175 edges, with normalized cut 0.104903.
They are connected parts of the large component, not separate components of
the original graph. Reversing the eigenvector sign exchanges the labels but
preserves membership.

\subsection*{Relationship to the original data}

The strongest raw-feature association is onset rate. Domain means account for
44.97\% of its total variance, measured by
\[
 \eta^2=\frac{\sum_g n_g(\bar x_g-\bar x)^2}
                  {\sum_i(x_i-\bar x)^2}.
\]
\begin{center}
\begin{tabular}{lrr}
\toprule
Measure & Domain 1 & Domain 2\\
\midrule
Mean onset rate (events/s) & 2.433 & 4.385\\
Mean BPM & 128.62 & 116.87\\
MFCC mean 01 (standardized mean) & 0.49 & $-0.46$\\
\bottomrule
\end{tabular}
\end{center}
Domain 2 has denser onsets but lower average BPM, so the partition is not simply
fast versus slow music. Several MFCC means also differ; for example, MFCC
mean 01 has $\eta^2=0.2264$, indicating a timbral contribution.

\begin{figure}[H]
\centering
\includegraphics[width=\linewidth]{latex_genres.pdf}
\caption{Prevalence of the 12 most common genre tags within each nodal domain.
Tracks may carry multiple labels.}
\label{fig:genres}
\end{figure}

Domain 1 has more classical, soundtrack, ambient and orchestral labels;
domain 2 has more electronic and pop labels. Classical prevalence is 23.16\%
versus 3.58\%, while electronic prevalence is 17.13\% versus 31.72\%.
The domains remain musically mixed. Acoustic associations are descriptive
because those features also define the graph; tags provide an interpretation,
not an independent classification test.

\section*{E7. Discussion}

The graph is almost connected at $k=2$ and fully connected at $k=3$, despite
low degrees and little triangle closure. Its spectral partition captures
differences in onset density, timbre and genre composition. Increasing $k$
strengthens normalized spectral connectivity, but does not necessarily improve
musical interpretation.

The two-stage sample, repeated artists and representation choices limit
generalization. High retained PCA variance does not guarantee preservation
of every nearest neighbor. Artist-aware resampling and parameter sensitivity
analysis would help assess robustness. These findings describe the specified
graph, rather than pure genre classes or predictive performance.

\paragraph{GitHub repository.}
\url{https://github.com/YiJin0716/MusicGraphAnalysis}

\begin{thebibliography}{9}
\bibitem{jamendo}
D. Bogdanov, M. Won, P. Tovstogan, A. Porter, and X. Serra.
\emph{The MTG-Jamendo Dataset for Automatic Music Tagging}.
Machine Learning for Music Discovery Workshop, ICML, 2019.
\url{https://hdl.handle.net/10230/42015}.
\bibitem{repository}
Music Technology Group. MTG-Jamendo dataset: metadata, features and documentation.
\url{https://github.com/MTG/mtg-jamendo-dataset}.
\end{thebibliography}
\end{document}
